\section{Experimental Design}
\label{sec:design}

\subsection{Four question shapes}

The benchmark has 20 questions in four categories of five, listed in full in
\Cref{app:questions}. Each category carries a pre-registered expectation of which
architecture should win, and two of the four were chosen so that the graph is the
\emph{wrong} tool:

\begin{itemize}
  \item \textbf{T1, single-document semantic.} What did one named company disclose
        about one topic in its risk factors? The whole answer lies inside one
        document. Expected winner: vector RAG or GraphRAG.
  \item \textbf{T2, structured aggregation.} Which audit firm serves the most
        companies; which jurisdiction of incorporation is most common. These are
        \texttt{GROUP BY} queries over typed columns. Expected winner:
        text-to-SQL, and a traversal should add nothing.
  \item \textbf{T3, multi-hop relational plus text.} For every company where a
        named person holds an insider role, list its subsidiaries and summarize
        its reported cybersecurity risks. Expected winner: GraphRAG.
  \item \textbf{T4, cross-document entity resolution.} Is a named subsidiary also
        an SEC filer in its own right, and which parent lists it? Expected
        winner: GraphRAG.
\end{itemize}

\subsection{Selecting subjects that make identity hard}

The T3 subjects were not chosen for convenience. Two conditions had to hold
simultaneously, both checked against the full-text index over all 5{,}210
risk-factor sections:

\begin{enumerate}
  \item the surname appears in \emph{zero} risk-factor sections, so no similarity
        search can reach the person at all; and
  \item two or more distinct filer identifiers share that surname, so resolving
        identity by name is demonstrably wrong rather than merely imprecise.
\end{enumerate}

Candidates were rejected on the first condition when the surname was too common in
the text (32, 49, 63 and 109 matching sections respectively) and on the second
when the surname turned out to be unique in the corpus. One unique-surname subject
was retained deliberately as a control: same question shape, no identity trap.

\subsection{Scoring}

Scoring is a deterministic function of the answer text. There is no LLM judge, for
the reason given in \Cref{sec:related}. An attempt \textbf{passes} when all four
conditions hold:

\begin{enumerate}
  \item every required entity is named (entity recall $=1.0$);
  \item at least one required topical keyword is present;
  \item no forbidden entity is named --- these are the specific wrong companies
        that a name-based rather than identifier-based resolution would return;
  \item where ground-truth citations exist, at least one is correct.
\end{enumerate}

The fourth condition is what makes provenance count. Without it, a model that
named the right companies but cited nothing would score identically to one that
cited the exact filings. Failures are then labelled: \textbf{falsehood} when the
answer is not correct, is not a refusal, and either names a forbidden entity or
cites a filing outside the valid set; \textbf{uncited} when the required entities
and keywords are present with nothing false but no correct citation;
\textbf{refused} when a regex over the answer detects an explicit statement that
the question cannot be answered from the evidence; and \textbf{fail} otherwise.
The labels partition the attempts, so the columns of \Cref{tab:main} sum to 60.

Aggregate questions have no ground-truth citation set. For those, any accession the
model cites is legitimate provenance and is not counted against it --- an earlier
version of the scorer did count it, and \Cref{sec:repro} records what that cost.

\subsection{Three attempts, all of them counted}

Every question was asked three times of every architecture, giving
$20 \times 3 \times 3 = 180$ scored attempts. The three attempts are
\emph{independent repeats to measure consistency, not retries}: every attempt is
scored and every attempt counts, and there is no best-of-three anywhere in the
reported numbers.

This is not ceremony. Even at temperature 0 the generations are not
deterministic, and \textbf{9 of the 60}
question-by-architecture cells returned different labels across their three
attempts --- 7 for text-to-SQL, 2 for GraphRAG and
0 for vector RAG. Text-to-SQL was inconsistent on four of the five
multi-hop questions, which is precisely where its false statements occur: an
evaluation drawing a single attempt per question could record an honest refusal
on the same question where another draw asserts something false.

We therefore report two metrics. \textbf{Attempts passed} out of 60 is the raw
rate. \textbf{Questions solved} out of 20 credits a question only when all three
of its attempts passed, and is the stricter reading: an answer that is correct
once in three is not a capability a practitioner can deploy. The two are not
related by a fixed ratio --- an architecture with $q$ questions solved has at
least $3q$ attempts passed, and the excess is contributed by questions solved
inconsistently.

\subsection{Fairness controls}

A rigged baseline would prove nothing, so we record exactly what was equalized:
identical model, temperature, reasoning setting and answer-token budget for all
three arms; an equal LLM call budget of two calls for the query-writing arms and
one for vector RAG, which has no query to write; the complete schema shown to
text-to-SQL so that every multi-hop answer is reachable by joins; the
surname-first naming convention documented in both schema descriptions; one
repair attempt for each query-writing arm; and the identical fan-out guardrail on
both query languages. Two of the four categories were selected so that the graph
should lose them.

We note one asymmetry we did not eliminate, because it favours the graph and we
would rather report it than hide it. The repair trigger fires on an exception or
zero rows. An over-broad SQL join returns \emph{many} rows and therefore never
trips it, even when the answer is wrong, whereas an exact graph property match
that misses returns nothing and earns a free second attempt. The same rule confers
unequal benefit; \Cref{tab:repair} quantifies the result.
